\needspace{10\baselineskip}
\hypertarget{native-322-zusatz17}{}
\pdfbookmark[1]{17. Zodiakus auf der Himmelskugel}{zusatz-017}
\begin{center}
[17.]

Zodiakus auf der Himmelskugel.

1847 Oct. 28.
\end{center}

Für jeden Punkt der Himmelskugel ist entweder gar keine, eine, oder mehr als eine Combination vorhanden, die die Sichtbarkeit daselbst des einen Planeten vom andern aus bedingt, oder eine bestimmte Zahl von Auflösungen. Beim Durchgange durch die Limitenlinie ändert sich allemal die Anzahl der Auflösungen um 2 Einheiten; so lange aber, bei irgend einem Wege auf der Himmelskugel, die Limitenlinie nicht getroffen wird, bleibt die Zahl der Auflösungen ungeändert. Die Frage ist nun, wie man die Seite der Limite, wo zwei Auflösungen mehr sind, von der andern, wo zwei weniger sind, unterscheidet. Folgende Betrachtungen dienen dazu.

1. Es sei \(p\) ein Punkt der Planetenbahn, von welcher aus die Punkte der andern Bahn auf die Himmelskugel projicirt werden; \(P\) der \SourceLeaf{323}correspondirende Punkt der andern Bahn; \(\Pi\) der Punkt der Knotenlinie, wo die Tangenten an den beiden Planetenbahnen in \(p\) und \(P\) einander schneiden; ferner sei
\[
p\Pi=n,\qquad P\Pi=N,\qquad pP=\Delta,
\]
\(S\) ein anderer beliebiger fester Punkt der Knotenlinie (etwa die Sonne). Die drei Figuren hat man sich in drei verschiedenen Ebenen vorzustellen: I in der Ebene durch \(\Pi pP\); II in der Ebene durch \(\Pi Sp\); III in der Ebene durch \(\Pi SP\). Noch bezeichnen wir mit
\begin{quote}
\noindent \(p,P\) die Winkel der Ebenen II und III mit I,\\
\(q,Q\) die Winkel \(p\Pi S,P\Pi S\),\\
\(u,U\) die Winkel \(\Pi pP,\Pi Pp\),\\
\(r,R\) die Krümmungshalbmesser der beiden Bahnen in \(p\) und \(P\).
\end{quote}

Es seien nun \(p^{\prime},P^{\prime}\) zwei andere Punkte correspondirend, und dazu der Punkt \(\Pi^{\prime}\) der Knotenlinie. \(\Pi\Pi^{\prime}=\delta\) ist als unendlich klein angenommen. Es werden dann die Normalen aus \(p^{\prime},P^{\prime}\) auf \(pP\):
\[
p^{*}p^{\prime}=\frac{\delta r\sin q\sin u}{n}\mathpunct{\raisebox{.45ex}{\normalfont,}}\qquad
P^{*}P^{\prime}=\frac{\delta R\sin Q\sin U}{N}\mathpunct{\cdot}
\]

Ferner seien \(p^{\prime\prime},P^{\prime\prime}\) zwei andere Punkte in den Tangenten \(p\Pi,P\Pi\), so dass \(p^{\prime\prime}P^{\prime\prime}\) parallel mit \(pP\), und \(\frac{pp^{\prime\prime}}{p\Pi}=\frac{PP^{\prime\prime}}{P\Pi}=\omega\) unendlich klein. Eine durch \(p^{\prime\prime}P^{\prime\prime}\) normal gegen I gelegte Ebene treffe die Curven in \(p^{\prime\prime\prime},P^{\prime\prime\prime}\), und \(p^{\mathrm{IV}},P^{\mathrm{IV}}\) seien die Fusspunkte der von \(p^{\prime\prime\prime},P^{\prime\prime\prime}\) auf I gefällten Perpendikel. Man hat dann
\[
\begin{aligned}
p^{\prime\prime\prime}p^{\mathrm{IV}}&=p^{\prime\prime}p^{\prime\prime\prime}\mathord{.}\sin p=\frac{nn\omega\omega}{2r}\sin p\\[5pt]
P^{\prime\prime\prime}P^{\mathrm{IV}}&=P^{\prime\prime}P^{\prime\prime\prime}\mathord{.}\sin P=\frac{NN\omega\omega}{2R}\sin P\\[5pt]
p^{\prime\prime\prime}p^{\mathrm{IV}}\mathord{.}p^{*}p^{\prime}&=\frac{\delta\omega\omega n\sin p\sin q\sin u}{2}\\[5pt]
P^{\prime\prime\prime}P^{\mathrm{IV}}\mathord{.}P^{*}P^{\prime}&=\frac{\delta\omega\omega N\sin P\sin Q\sin U}{2}\mathpunct{\cdot}
\end{aligned}
\]

Man sieht aber leicht, dass \(n\sin u=N\sin U\) das Perpendikel aus \(\Pi\) auf \(pP\); ferner \(\sin p\sin q=\sin P\sin Q={}\) Sinus des Winkels, welchen die Linie \(\Pi S\) mit der Ebene I, und folglich mit dem gedachten Perpendikel macht, woraus man leicht schliesst, dass \(n\sin p\sin q\sin u=N\sin P\sin Q\sin U\) nichts anderes ist, als die kürzeste Entfernung der Geraden \(pP,\Pi S\) von einander.

\SourceLeaf{324}Endresultat ist also
\[
p^{\prime\prime\prime}p^{\mathrm{IV}}\mathord{.}p^{*}p^{\prime}=P^{\prime\prime\prime}P^{\mathrm{IV}}\mathord{.}P^{*}P^{\prime}.
\]
Ist folglich \(p^{\prime\prime\prime}p^{\mathrm{IV}}>P^{\prime\prime\prime}P^{\mathrm{IV}}\), so ist \(p^{*}p^{\prime}<P^{*}P^{\prime}\) und umgekehrt. Die erste Ungleichheit zeigt aber an, dass eine mit \(p^{\prime\prime\prime}P^{\prime\prime\prime}\) parallel aus \(p\) gezogene Gerade einerseits und \(S\) andererseits auf entgegengesetzte Seite von I fallen; die zweite, dass die Projectionen auf die Himmelskugel von den beiden Geraden \(p^{\prime}P^{\prime}\) und \(pP^{\prime}\) (wovon die letztere nur um ein unendlich kleines 2\textsuperscript{ter} Ordnung von der Limitenlinie absteht) auf Einer Seite von der Projection von \(pP\) liegen.

Hieraus ist das gewünschte Criterium leicht abzuleiten; man kann es so vortragen:

[Sei] \(B\) [die] äussere Planetenbahn von der Sonne aus auf die Himmelskugel bezogen; dieser grösste Kreis theilt die Himmelskugel in zwei Hälften (1) und (2). Es sei (1) diejenige, in welcher der heliocentrische Ort des innern Planeten erscheint. Dann liegt gegen ein Stück der Limite \(LL^{\prime}\) diejenige Gegend, wo zwei Auflösungen mehr gelten, eben so wie (1) gegen \(B\), indem man in \(LL^{\prime}\) und \(B\) als vorwärts gehenden Sinn Zusammengehöriges annimmt. Es ist hiebei angenommen, dass der Planet der äussern Bahn der beobachtete ist. Im entgegengesetzten Fall ist es umgekehrt.

Symbolisch und noch allgemeiner kann das Criterium so ausgedrückt werden:

Es seien \(l,l^{\prime}\) zwei einander unendlich nahe Punkte in einer Limite; \(p,p^{\prime}\) die entsprechenden heliocentrischen Örter des Planeten, von welchem aus [beobachtet wird], \(P,P^{\prime}\) hingegen die des Planeten, welcher beobachtet wird; \(a\) ein Punkt nahe an \(ll^{\prime}\) auf der Seite, wo zwei Auflösungen mehr gelten als auf der andern. Dann ist, wenn das Congruenzzeichen \(\equiv\) zur Bezeichnung gleicher Lageordnung [dient], und ein vorgesetztes \(-\)Zeichen die entgegengesetzte Lageordnung bedeutet:
\[
ll^{\prime}a\equiv PP^{\prime}p\equiv-pp^{\prime}P.
\]
\begin{center}\rule{.22\textwidth}{.3pt}\end{center}
